\answer{ex:09:1}{(a)~$P_\infty=0.2$，记忆$20\,\text{s}$。(b)~$P_\infty=1$，记忆$1\,\text{s}$。(c)~记忆与噪声幅度成正比：六倍。}
\answer{ex:09:2}{$P(t)=P_\infty\coth(t\sqrt{q/R})$。(a)~$\frac12\sqrt{R/q}=5\,\text{s}$——比记忆短一半。(b)~$t=\frac12\ln21\,\sqrt{R/q}\approx15\,\text{s}$：升高的\nt{ACET}应当维持这么久。}
\answer{ex:09:3}{方差为$qT/2+R/(2T)$；$T^*=\sqrt{R/q}$，方差为$\sqrt{qR}=P_\infty$，与滤波器~\eqref{eq:kalman}相同；增大$(\kappa+1/\kappa)/2$倍：$1.25$和$5.05$。}
\answer{ex:09:4}{$a>a_w=1-\theta/(mw)$：$w=0.041$时为$0.15$，$w=0.045$时为$0.22$。}
\answer{ex:09:5}{在$M\approx40$--$60$时开始出错，$M=80$时多出$1.9$个神经元，$M=100$时模式比例为$0.86$；来自其他模式的干扰方差$\approx(M-1)p(1-p)/N$，所以$M_{\max}\approx80$，$M_{\max}\propto N/p$。}
\answer{ex:09:6}{例如$\eta(a)=\eta\min\bigl(1,\max(0,(a-0.3)/0.6)\bigr)$。用$\eta a$时三个外来神经元全部进入模式；用$\eta(a)$时一个也没有，$a\ge0.9$时的写入不变（比例$1.00$）。}
\answer{ex:09:7}{(a)~$40$次回放；用习题~\ref{ex:09:6}的$\eta(a)$则永远不够。(b)~$200$次回放，$10$个夜晚。}
